\hmsubsection{Development Phases}{}{}\label{sec:development-phases}
We organised the project into four phases. The lengths of the phases
are not equal. Phase 3 took the most time because that is where
hardware and software had to be made to work together.
\subsubsection*{Phase 1 -- Inception (weeks 1--3)}
We met the supervisor and the partner and clarified the assignment.
We agreed on roles (Section \ref{sec:team}). We read  Sortify \cite{sortify} as the closest prior
project. We agreed on the minimum viable product
(Section \ref{sec:mvp}). At the end of the phase we had a first draft
of the requirements list and the risk register.
\subsubsection*{Phase 2 -- Bring-up (weeks 4--8)}
We ordered the hardware and brought it up one component at a time.
The OpenCM 9.04 controller was tested first and rejected because the
board we got stopped responding after we soldered the headers
(Section \ref{sec:controller-selection}). The OpenRB-150 was ordered
next and worked on the first try. We set up the OAK-D S2 alongside it.
A first version of the vision pipeline used HSV thresholding alone.
A first version of the firmware could drive each servo on its own.
At the end of the phase the camera and the arm were both working, but
not yet talking to each other.
\subsubsection*{Phase 3 -- Integration (weeks 9--16)}
This is where most of the work happened. The two halves were joined,
and the Pi started sending real commands to the arm. We tried
MobileNetV2 during this phase and dropped it after measuring 200--500
ms per frame (Section \ref{sec:related-detection}). We added the
Hough Circle Transform alongside HSV to form the ensemble
(Section \ref{sec:ensemble}). The inverse kinematics solver was
written and tested against the physical arm. The sag model was
calibrated for the first time (Section \ref{sec:sag}).
\subsubsection*{Phase 4 -- Validation (weeks 17--20)}
The last phase was for testing and writing. We ran the system end to
end against the acceptance criteria and measured throughput and
reliability. The touch calibration procedure
(Section \ref{sec:touch-cal}) replaced the earlier ruler-based one.
The SVM colour verifier was disabled after it kept misclassifying dark
navy-blue balls (Section \ref{sec:svm}). We finalised the report and
prepared the demonstrator for the bachelor defence.
The phases match the sprint structure described in
Section \ref{sec:methodology}. Each phase contained two to four
sprints, and the sprint retrospectives are the natural place where we
moved from one phase to the next.
